\documentclass[11pt]{article}

\usepackage[margin=1in]{geometry}
\usepackage{xcolor}
\usepackage{titlesec}
\usepackage{tcolorbox}
\usepackage{amsmath,amssymb}
\usepackage{hyperref}
\usepackage{enumitem}
\usepackage{fancyhdr}
\usepackage{lmodern}

\tcbuselibrary{skins}

\definecolor{mauve}{rgb}{0.58,0,0.82}
\definecolor{headerblue}{rgb}{0.85,0.87,0.98}
\definecolor{deliverblue}{rgb}{0.06,0.10,0.45}
\definecolor{lightbg}{rgb}{0.96,0.97,1.0}
\definecolor{tabgold}{rgb}{0.99,0.95,0.80}

\hypersetup{colorlinks=true, linkcolor=blue, urlcolor=blue, citecolor=blue}

% "Problem N: Title" banner
\newcommand{\problem}[1]{%
  \vspace{1em}
  \noindent\colorbox{headerblue}{\parbox{\dimexpr\textwidth-2\fboxsep}{%
    \vspace{2pt}\textbf{\large #1}\vspace{2pt}}}
  \vspace{0.6em}
}

% [Question X.Y] (Kind, N points)
\newcommand{\question}[3]{%
  \vspace{0.9em}
  \noindent\textbf{[Question #1]}\ \ (\textit{\textbf{#2, #3 points}})\quad
}

% Boxed deliverables callout with a gold tab
\newtcolorbox{deliverables}[1]{
  enhanced, colback=lightbg, colframe=deliverblue, boxrule=0.9pt,
  arc=2pt, left=8pt, right=8pt, top=12pt, bottom=8pt,
  overlay={\node[fill=tabgold, draw=deliverblue, line width=0.9pt,
    rounded corners=1pt, inner xsep=8pt, inner ysep=3pt,
    font=\bfseries\footnotesize, anchor=north]
    at ([yshift=6pt]frame.north) {DELIVERABLES FOR #1};},
  before skip=14pt, after skip=10pt
}

\pagestyle{fancy}
\fancyhf{}
\lhead{\footnotesize 11-711 Fall 2026; Homework 1}
\rhead{\footnotesize\thepage}
\renewcommand{\headrulewidth}{0.4pt}

% Long \texttt{} paths cannot hyphenate, and the \question header glues itself to
% the text that follows, which pushed nine paragraphs past the right margin.
% emergencystretch lets TeX stretch interword space as a last resort rather than
% overflowing; \sloppy prefers loose lines to overfull ones.
\emergencystretch=3em
\sloppy
\hbadness=10000

\begin{document}

\begin{center}
  {\Huge\sffamily\color{mauve} Homework 1}\\[6pt]
  {\Large\sffamily Build a Language Model}\\[10pt]
  {\large 11-711 Advanced Natural Language Processing, Fall 2026}\\[4pt]
  \textbf{Released:} September 1 \quad\textbf{Due:} September 24, 11:59 PM\\[2pt]
  \textbf{Late submissions close:} October 14, 11:59 PM\\[2pt]
  \textbf{Individual assignment} \quad 200 points \quad 10\% of your grade
\end{center}

\vspace{0.8em}
\hrule
\vspace{0.8em}

\noindent In this homework you will build a language model end to end. You will train a
tokenizer, implement the Llama-2 architecture, pretrain a model under a fixed compute budget,
evaluate what it learned, and adapt it to a downstream task.

\vspace{0.6em}
\noindent\textcolor{red}{\textbf{There are two Gradescope submissions.} Upload
\texttt{submission.zip} (built by \texttt{scripts/zip-submission.sh}) to \textbf{Assignment1: Code}
and \texttt{report.pdf} to \textbf{Assignment1: Report}. Coding questions are graded only by the
autograder, written questions only from the report; nothing is submitted twice. Use the provided
LaTeX template, and when uploading the report, mark the pages that answer each question.}

\vspace{0.6em}
\noindent\textbf{Late policy.} Per the syllabus, each day (or part of a day) a submission is late
costs 5\% of your grade for this homework; both submissions must be in by October 14, after which
Gradescope closes. Requests for excused late days go to the course email at least 7 days before the
deadline.

\vspace{0.6em}
\noindent\textbf{On the use of AI.} Per the syllabus, you may use AI to help you write \emph{code}.
You may \textbf{not} use it to write any part of your report. Question 6 asks you to describe how
you used AI; per the syllabus, using AI without making an honest
attempt at that question is itself an academic integrity violation.

\vspace{0.6em}
\noindent\textbf{Roughly half the points are written analysis}, and those questions ask about
\emph{your} runs: your loss curves, your search table, your generations. Numbers that do not
match your submitted checkpoint will not receive credit.

%======================================================================
\problem{Problem 1: Tokenization (35 points)}

\noindent A tokenizer maps a string $s \in \mathcal{A}^*$ over an alphabet $\mathcal{A}$ to a sequence of
tokens $\mathbf{x} \in \mathcal{V}^*$. Tokenizers are complete in their alphabet: any string over
$\mathcal{A}$ can be tokenized. Everything downstream, what the model can represent, how much
context a document consumes, how much it costs to serve, inherits from this choice.

\question{1.1}{Writing}{12}
We might reasonably expect a tokenizer to satisfy the following:

\begin{itemize}[nosep, leftmargin=2em]
  \item \textbf{Injective}: for all $s \neq s'$, $\mathrm{tokenize}(s) \neq \mathrm{tokenize}(s')$.
  \item \textbf{Invertible}: for any string $s$, $s = \mathrm{detokenize}(\mathrm{tokenize}(s))$.
  \item \textbf{Preserves concatenation}: $\mathrm{tokenize}(s + s') = \mathrm{tokenize}(s) + \mathrm{tokenize}(s')$.
\end{itemize}

All three are violated by the tokenizers of widely used language models.

\begin{deliverables}{Q1.1}
\begin{enumerate}[label=\Alph*., leftmargin=1.6em]
  \item Find counterexamples showing that a popular, publicly available tokenizer fails each
    property. \textbf{Use \texttt{google-bert/bert-base-cased}.} The autograder has it
    cached and may have no network access, so another choice risks failing there even if it works
    locally. Write code in \texttt{tests/test\_tokenizer.py} so that the assertions in
    \texttt{test\_not\_injective()}, \texttt{test\_not\_invertible()} and
    \texttt{test\_not\_preserving\_concat()} all pass. In your report, describe in at most two
    sentences per test what violation you found.
  \item Explain why it is essentially impossible to build a non-trivial tokenizer that preserves
    concatenation.\footnote{A trivial tokenizer is one whose vocabulary is the alphabet itself.}
\end{enumerate}
\end{deliverables}

\question{1.2}{Coding}{15}
Byte-pair encoding (BPE) \emph{learns} a tokenizer from a corpus. It starts with the alphabet as
its vocabulary and repeatedly merges the most frequent adjacent pair of tokens into a new token.
Each iteration adds one token while keeping every existing one.

\begin{deliverables}{Q1.2}
\begin{enumerate}[label=\Alph*., leftmargin=1.6em]
  \item Implement \texttt{ASCIIBPETokenizer.merge()} in \texttt{src/tokenizer/bpe.py}. It must
    pass \texttt{test\_one\_merge()}, \texttt{test\_two\_merges()} and
    \texttt{test\_100\_merges()} in \texttt{tests/test\_tokenizer.py}.
    \emph{You are given helpers} in the same file: \texttt{string\_to\_ascii()},
    \texttt{compute\_bigram\_statistics()} and \texttt{replace\_bigram()}. Read them first, you
    should not need to write your own.
  \item Implement \texttt{encode()} and \texttt{decode()}. These must pass \texttt{test\_encode()}
    and \texttt{test\_decode()}.
\end{enumerate}
Additional hidden tests may be run after the deadline.
\end{deliverables}

\question{1.3}{Writing}{4}
We have provided a Unicode BPE tokenizer trained only on English text
(\texttt{data/english-tokenizer.json}). Run \texttt{python -m src.tokenizer.visualize} to load it
and inspect the tokens it learned.

\begin{deliverables}{Q1.3}
Answer in at most two sentences each.
\begin{enumerate}[label=\Alph*., leftmargin=1.6em]
  \item What is the longest token you see, and what does that tell you about the training data?
  \item How can BPE training compromise the privacy of the corpus it was trained on?
\end{enumerate}
\end{deliverables}

\question{1.4}{Writing}{4}
A tokenizer trained on one distribution spends more tokens per unit of meaning on anything else.
Using the Llama SentencePiece tokenizer (\texttt{src/llama/sp\_tokenizer.py}), measure how many
tokens it takes to encode comparable amounts of content across different kinds of text.
\texttt{python -m src.tokenizer.inspect\_tokenizer} runs the core measurements:
it prints how sample digits tokenize and tokens per character for sample texts in English, source
code and a non-Latin script; \texttt{-\/-text} covers anything else, including the exact number
ranges in A.

\begin{deliverables}{Q1.4}
Answer in at most three sentences each.
\begin{enumerate}[label=\Alph*., leftmargin=1.6em]
  \item \textbf{Digits.} Tokenize the numbers 1--20, then 1000--1020. How are multi-digit numbers
    split? Does the same number of digits always produce the same number of tokens? Explain what
    this tokenization choice means for a model learning arithmetic, compared to a tokenizer that
    merges frequent digit sequences into single tokens.
  \item \textbf{Domains.} Compare tokens-per-character for three texts of similar length: English
    prose, source code, and a passage in a non-Latin script. Which is most expensive, and why?
  \item \textbf{Who pays.} API pricing and context windows are both denominated in tokens. State
    one concrete consequence of your answer to (B) for users writing in the most expensive of the
    three.
\end{enumerate}
\end{deliverables}

%======================================================================
\problem{Problem 2: Implementing Llama-2 (37 points)}

\noindent You will implement the core of the Llama-2 architecture. All relevant code is in
\texttt{src/llama/}. \textbf{Do not modify the classes, functions or their
arguments}, and do not import new dependencies; both will break the autograder. Aim for an
efficient implementation: call PyTorch operations over whole tensors rather than looping. You may
not use layers that make the exercise trivial (e.g.\
\texttt{torch.nn.TransformerDecoder} or \texttt{F.scaled\_dot\_product\_attention}); if you are
unsure whether something is allowed, ask on Piazza.

\question{2.1}{Coding}{24}
Implement the following. Each is marked \texttt{\# todo} in the source.

\begin{itemize}[nosep, leftmargin=2em]
  \item \texttt{RMSNorm.\_norm}: root-mean-square normalization.
  \item \texttt{Attention.compute\_query\_key\_value\_scores}: scaled dot-product attention with
    a causal mask.
  \item \texttt{LlamaLayer.forward}: one transformer block.
  \item \texttt{Llama.generate}: temperature and top-$p$ (nucleus) sampling.
  \item \texttt{apply\_rotary\_emb} in \texttt{src/llama/rope.py}: rotary position embeddings.
    This one is fiddly; \texttt{tests/test\_llama.py} tests it in isolation.
\end{itemize}

\noindent If you get stuck, \href{https://arxiv.org/pdf/1706.03762}{Attention Is All You Need}
(Vaswani et al., 2017) and \href{https://nlp.seas.harvard.edu/annotated-transformer/}{The Annotated
Transformer} are useful references for understanding the mechanics of attention and transformer
blocks in code.

\begin{deliverables}{Q2.1}
Upload your code to Gradescope and make sure the tests in \texttt{tests/test\_llama.py} pass.

\vspace{4pt}
\noindent The test that matters most is \texttt{test\_llama\_forward\_matches\_reference}. It loads
real pretrained weights (\texttt{stories42M.pt}, an 8-layer, 42M-parameter Llama trained on
TinyStories) and checks your model reproduces their activations. Shape-correct but subtly wrong
code fails it. If the reported maximum difference is around $10^{0}$, suspect your normalization or
attention; around $10^{-4}$ is ordinary floating-point noise.
\end{deliverables}

\question{2.2}{Writing}{6}
\href{https://arxiv.org/abs/1608.05859}{Press and Wolf (2017)} propose tying the input embedding
and the output projection.

\begin{deliverables}{Q2.2}
\begin{enumerate}[label=\Alph*., leftmargin=1.6em]
  \item In at most three sentences, explain what weight tying does and why it helps.
  \item Let $d$ be the hidden size, $v$ the vocabulary size, $b$ the batch size and $s$ the
    sequence length. Given hidden states $h \in \mathbb{R}^{b \times s \times d}$ and token
    embeddings $E \in \mathbb{R}^{v \times d}$, write one line of PyTorch that computes the token
    logits using weight tying.
\end{enumerate}
\end{deliverables}

\question{2.3}{Writing}{7}
Llama-2 uses grouped-query attention: several query heads share one key/value head. During
generation the model caches keys and values for every position it has already seen.

\begin{deliverables}{Q2.3}
\begin{enumerate}[label=\Alph*., leftmargin=1.6em]
  \item Write an expression for the number of bytes the KV cache occupies, in terms of the number
    of layers $L$, key/value heads $n_{kv}$, head dimension $d_h$, sequence length $s$, batch size
    $b$, and bytes per element.
  \item Using the \texttt{configs/llama-base.yaml} values, compute the cache size for batch
    size $1$ at $s = 256$ and, hypothetically, at $s = 8192$, in \texttt{float16}.
  \item GQA reduces $n_{kv}$ while leaving the number of query heads unchanged. Given your
    expression, what does that buy, and what is the cost? Why is this a bigger deal at inference
    time than during training?
\end{enumerate}
\end{deliverables}


%======================================================================
\problem{Problem 3: Pretraining under a compute budget (68 points)}

\noindent You will now pretrain your model on \href{https://arxiv.org/abs/2305.07759}{TinyStories}, a corpus
of machine-generated children's stories. We provide the corpus pre-tokenized, so you do not need to
run your Part 1 tokenizer over it.
Training is budgeted in \textbf{FLOPs rather than wall-clock time}, so the assignment is fair
across hardware. \texttt{src/train.py} prints the projected cost before training begins and warns
if you exceed the final $10^{17}$ budget; for Q3.5's per-configuration budget, check the printed
projection against $10^{15}$ yourself. Total cost is
\begin{align*}
  \text{FLOPs} = \ &\texttt{flops\_per\_token} \times \texttt{num\_training\_steps} \\
  \times\ &\texttt{grad\_accumulation\_steps} \times \texttt{batch\_size} \times \texttt{seq\_len},
\end{align*}
where $\texttt{flops\_per\_token} = 6N + 12 L d\, s_{\max}$ for a model with $N$ parameters, $L$
layers and hidden size $d$. Every hyperparameter moves this: a larger model buys you fewer steps.

\question{3.1}{Writing}{5}
\texttt{tests/test\_train.py} contains \texttt{test\_train\_sanity\_check}, which trains on a
sequence of uniformly random tokens drawn from $\{0,\dots,7\}$ and \emph{asserts the loss stays
above} $\log(8) \approx 2.08$.

\begin{deliverables}{Q3.1}
\begin{enumerate}[label=\Alph*., leftmargin=1.6em]
  \item Why is $\log(8)$ the floor for this data? Show the reasoning.
  \item Describe a specific bug in the attention implementation that would let a model score
    \emph{below} this floor, and explain the mechanism.
  \item This test is worth zero points but failing it invalidates Questions 3.4--3.6. Why would
    those results be meaningless?
\end{enumerate}
\end{deliverables}

\question{3.2}{Coding}{28}
Implement the following in \texttt{src/train.py}:
\texttt{random\_batch\_sampler}, \texttt{sequential\_batch\_sampler}, \texttt{cosine\_lr\_schedule}
and \texttt{compute\_language\_modeling\_loss}.

Also implement \texttt{AdamW.step} in \texttt{src/optim/adamw.py}, with decoupled weight decay and
gradient clipping. Use the efficient bias correction from the end of Section 2 of Kingma \& Ba
(2014) rather than forming $\hat{m}$ and $\hat{v}$, and fold the learning rate into the weight decay
update. Your training runs in Part 3 use \emph{your} optimizer, not \texttt{torch.optim.AdamW}.

\begin{deliverables}{Q3.2}
Upload your code and make sure \texttt{tests/test\_train.py} and
\texttt{tests/test\_optimizer.py} pass, including \texttt{test\_train\_sanity\_check} (see Q3.1).
\end{deliverables}

\question{3.3}{Writing}{6}
Cosine annealing with warmup has two phases. During warmup, $t \in [0, a)$, the learning rate rises
linearly from $0$ to $\mathrm{lr}_{\max}$. During annealing, $t \in [a, b)$, it decays from
$\mathrm{lr}_{\max}$ to $\mathrm{lr}_{\min}$ along a half-cosine. For $t \geq b$ it stays at
$\mathrm{lr}_{\min}$.

\begin{deliverables}{Q3.3}
\begin{enumerate}[label=\Alph*., leftmargin=1.6em]
  \item Write the expression for the learning rate during warmup, and during annealing.
  \item In at most two sentences: why warm up at all, and what does annealing buy over a constant
    learning rate?
\end{enumerate}
\end{deliverables}

\question{3.4}{Writing}{5}
Train the smoke-test configuration: \texttt{python -m src.train configs/llama-tiny.yaml}. It takes
about a minute and is not meant to produce a good model.

\begin{deliverables}{Q3.4}
\begin{enumerate}[label=\Alph*., leftmargin=1.6em]
  \item Report the final validation loss.
  \item Roughly how many steps before the training loss stops improving?
  \item Include the training curve (a screenshot from Weights \& Biases is fine).
\end{enumerate}
\end{deliverables}

\question{3.5}{Writing}{14}
Now search for a good configuration. \textbf{Each configuration you try must use at most $10^{15}$
FLOPs} --- about four minutes each on a Google Colab T4 instance. Within that budget you may change anything
that affects cost: \texttt{dim}, \texttt{n\_layers}, \texttt{n\_heads}, \texttt{n\_kv\_heads},
\texttt{max\_seq\_len}, \texttt{batch\_size}, \texttt{seq\_len},
\texttt{grad\_accumulation\_steps}, \texttt{num\_training\_steps}, and the learning rates. Evaluate
using validation perplexity.

\begin{deliverables}{Q3.5}
Report \textbf{at most three} configurations. For each, state what you changed and the validation
perplexity you got.
\begin{enumerate}[label=\Alph*., leftmargin=1.6em]
  \item Describe your experimental procedure in one paragraph. What did you vary, and why those
    knobs rather than others?
  \item Describe your results in one paragraph. Which change mattered most?
  \item Paste the YAML of your best configuration.
  \item You held compute fixed and traded model size against training steps. What does your data
    suggest about that trade-off, and how does it relate to what you know about scaling laws?
\end{enumerate}
\end{deliverables}

\question{3.6}{Coding}{10}
Train your final model using the configuration you chose. You may spend up to $10^{17}$ FLOPs,
100$\times$ the search budget, so budget roughly six hours on a Google Colab T4 instance.

\textbf{Keep the model architecture you reported in Q3.5}: \texttt{dim}, \texttt{n\_layers},
\texttt{n\_heads}, \texttt{n\_kv\_heads} and \texttt{max\_seq\_len}. Everything else is yours to
change in order to use the larger budget: \texttt{num\_training\_steps},
\texttt{grad\_accumulation\_steps}, \texttt{batch\_size}, \texttt{seq\_len} (up to
\texttt{max\_seq\_len}), and the learning rates. In practice you will mostly be raising
\texttt{num\_training\_steps}.

\begin{deliverables}{Q3.6}
Report your final validation loss and perplexity, and submit \texttt{model.pt} and
\texttt{config.yaml} as instructed in the README. A hidden test measures perplexity on a held-out
split you have not seen: \textbf{perplexity $\leq 5.5$ earns full credit, $\leq 8.0$ earns half}.

\textbf{Your own validation perplexity is a good proxy for this}, so you do not need to guess. The
held-out split is drawn from the same distribution as the validation set you can see, and in our
experience the two agree closely. If your validation perplexity is comfortably under 5.5, you are
done and can stop tuning. For scale, the smoke-test configuration \texttt{llama-tiny} reaches
roughly 128. Gradescope caps uploads at 100MB; \texttt{train.py} warns you if your checkpoint is too
large.
\end{deliverables}

%======================================================================
\problem{Problem 4: Generation and perplexity (35 points)}

\question{4.1}{Coding}{12}
Implement \texttt{softmax\_with\_temperature} in \texttt{src/generate.py}, and the top-$p$ sampling
branch of \texttt{Llama.generate} if you have not already (Q2.1).

\begin{deliverables}{Q4.1}
Upload your code and make sure \texttt{tests/test\_generate.py} passes.
\end{deliverables}

\question{4.2}{Writing}{11}
Compare your Part 3 model against \texttt{stories42M}. Both are Llamas trained on TinyStories, so
the architecture family is held fixed; what differs is model size and how much training went in. Choose
three prefixes, from \texttt{data/prefixes.jsonl} or your own, and generate at least 32 tokens
from each model, at several temperatures. \texttt{python -m src.compare\_generations
-\/-config <your config> -\/-checkpoint <your model.pt>} does this side by side for both models
(temperatures 0.0, 0.7, 1.5 by default) and writes every generation to a jsonl file.

\begin{deliverables}{Q4.2}
\begin{enumerate}[label=\Alph*., leftmargin=1.6em]
  \item Report your prefixes, the temperatures and \texttt{max\_new\_tokens} you used, and paste
    the generations from each model separately.
  \item What happens to the generations as temperature approaches zero, and as it grows large?
    Answer in at most two sentences.
  \item At a fixed temperature, compare the quality and perplexity of the two models' output.
    Report what you find and, in at most three sentences, why.
  \item Your model and \texttt{stories42M} differ in both size and training tokens. What can you
    \emph{not} conclude from this comparison, and what experiment would separate the two?
\end{enumerate}
\end{deliverables}

\question{4.3}{Writing}{12}
Now measure perplexity directly. \texttt{src/evaluate\_perplexity.py} scores every document in a
jsonl file and reports perplexity (its docstring has more detail):

\begin{quote}\small\ttfamily
python -m src.evaluate\_perplexity -\/-checkpoint data/stories42M.pt \\
\hspace*{2em}-\/-documents data/ppl\_examples.jsonl
\end{quote}

\noindent \texttt{data/ppl\_examples.jsonl} already contains an in-domain TinyStories passage and a
Wikipedia passage. Pass \texttt{-\/-strip-word the} to delete a word before scoring, and
\texttt{-\/-config <your config> -\/-checkpoint <your model.pt>} to score with your own Part 3 model
instead.

\begin{deliverables}{Q4.3}
\begin{enumerate}[label=\Alph*., leftmargin=1.6em]
  \item \textbf{Function words.} Compute your model's perplexity on a held-out TinyStories sample.
    Then delete every instance of ``the'' and recompute. Report both numbers and explain the change
    in at most three sentences.
  \item \textbf{Distribution shift.} Compute your model's perplexity on a few paragraphs of English
    Wikipedia. Compare it to (A). Explain the gap, and state what that implies about reporting a
    single perplexity number for a model.

\end{enumerate}
\end{deliverables}

%======================================================================
\newpage

\problem{Problem 5: Adapting the model (25 points)}

\noindent Pretraining gives you a model that continues text. Most of what people want from language models
is something else. Here you adapt the pretrained model to sentiment classification, on SST-5
(5-way) and CFIMDB (binary).

\question{5.1}{Coding}{10}
Implement \texttt{LlamaEmbeddingClassifier.forward} in \texttt{src/adapt.py}: pool a single vector
out of the sequence, apply dropout, pass it through the classifier head, and return
log-probabilities.

\begin{deliverables}{Q5.1}
Upload your code and make sure \texttt{tests/test\_adapt.py} passes.

\vspace{4pt}
\noindent Note the padding-invariance tests. Batches are right-padded to their longest sequence, so
\texttt{hidden\_states[:, -1]} is a \emph{padding} position for every sequence that is not the
longest in its batch. Your pooling must find the last real token.
\end{deliverables}

\question{5.2}{Writing}{8}
Run the model in two settings against \texttt{stories42M}: zero-shot prompting
(\texttt{--option prompt}) and finetuning with a classification head
(\texttt{--option finetune}), on both datasets.

\begin{deliverables}{Q5.2}
\begin{enumerate}[label=\Alph*., leftmargin=1.6em]
  \item Report dev accuracy for all four runs (two settings $\times$ two datasets). Report dev
    only: your copy of \texttt{cfimdb-test.txt} carries placeholder labels (every label is 0),
    so a locally computed test accuracy is meaningless there, and your runs print \texttt{n/a}
    for it. Every run writes its prediction files to \texttt{outputs/preds/}; they are included
    in your code submission, and we score the test split against the real held-out labels
    ourselves.
  \item Chance is 0.200 for 5-way SST-5 and 0.500 for binary CFIMDB, but chance is only an
    honest floor if a constant predictor cannot beat it. Compute the majority-class baseline of
    each dev set: the accuracy of always predicting its most common label. Then look at what
    your prompting run actually predicts. The first field of each line of the dev output file
    is the predicted label, so \texttt{cut -d'|' -f1 <dev-out> | sort | uniq -c} prints the
    whole distribution. Report both baselines and both distributions, and say in at most four
    sentences what they imply about whether zero-shot prompting is doing sentiment
    classification here. Give one reason about the model and one about the single-token
    scoring setup.
  \item Finetuning does much better while starting from the same weights. What has changed, given
    that the pretrained model has seen no sentiment labels at all?
\end{enumerate}
\end{deliverables}

\question{5.3}{Writing}{7}
Now finetune \emph{your own} Part 3 model on SST-5, using the same procedure and passing
\texttt{-\/-backbone-config} with the config your model was trained with. (SST-5 only, to keep
the runs short.)

\begin{deliverables}{Q5.3}
\begin{enumerate}[label=\Alph*., leftmargin=1.6em]
  \item Report dev accuracy for your model and for \texttt{stories42M}, finetuned identically.
  \item Is the difference between the two larger than run-to-run variation? Finetune at least twice
    with different \texttt{--seed} values to find out. Then, in at most three sentences, say what
    your answer implies about what the extra pretraining compute behind \texttt{stories42M} bought
    on this task. \textbf{``The difference is within noise'' is a perfectly good answer} if that is
    what you measure.
  \item Both models were pretrained only on children's stories, yet both can be finetuned for movie
    review sentiment. What does that tell you about what pretraining actually provides?
\end{enumerate}
\end{deliverables}

%======================================================================
\problem{Question 6: Use of AI (required, ungraded)}

\begin{deliverables}{Q6}
Describe how you used AI tools on this assignment: which parts, and how. If you did not use any,
say so. Per the syllabus this question is required, and failing to make an honest attempt at it is
itself an academic integrity violation.
\end{deliverables}

%======================================================================
\problem{Bonus (up to 10 points)}

\noindent Optional. Pick \textbf{one} and write it up in an appendix to your report: what you did, what you
measured, and what you concluded. Include plots and training logs.

\begin{itemize}[leftmargin=2em]
  \item \textbf{The tiniest adding Llama.} Pretrain a model to do 4-digit addition
    (\texttt{bonus/}). Find the smallest architecture reaching 100\% test accuracy within the
    \texttt{-\/-capacity} budget of \texttt{bonus/addition\_run.py}; the todos live in
    \texttt{bonus/addition\_data\_generation.py} and \texttt{bonus/addition\_run.py}. Report at least five ablations, including at least two that fail.
  \item \textbf{LoRA.} Implement \texttt{LoRALayer.forward} in \texttt{bonus/lora.py} and repeat
    Q5.2 with it. Report the fraction of parameters you trained (see
    \texttt{count\_lora\_parameters}) alongside the accuracy, and compare against full finetuning.
    Note that the update is scaled by $\alpha/r$; explain what you chose and why.
  \item \textbf{Something else.} Continued pretraining for domain adaptation; an alternative
    positional encoding; a different attention variant; alternative decoding strategies. Clear it
    on Piazza first.
\end{itemize}

%======================================================================
\vspace{1.2em}
\hrule
\vspace{0.8em}

\noindent\textbf{Acknowledgements.} This code is based on llama2.c by Andrej Karpathy (MIT, see
\texttt{LICENSE}). Parts of the code are also from the \texttt{transformers} library (Apache
License 2.0), and it contains adaptations from nanoGPT and PyTorch (see \texttt{copyright/}).

\end{document}
